\begin{abstract}
Agents that revisit evolving data can repeat the same procedural mistake even as the values change: using evidence released after a cutoff, confusing publication time with reporting time, or failing to ground a numerical change in contemporaneous documents. We study an identification problem rather than a skill-learning problem: when does inserting a reusable temporal procedure actually repair such a recurring failure, and what part of the intervention deserves credit? Our base design compares a targeted procedure (T), a complexity-matched target-blind helper (G), and the original agent (N); a binding repair requires T to beat both G and N. This extra anchor changes an observed verdict: on one DeepSeek cutoff cell, T=100\% and G=72\%, but N=100\%, so the apparent +28-point two-arm gain is not a repair. We then strengthen identification with two prospectively frozen controls. A no-op helper G$_0$ is behaviorally neutral relative to fresh N on independent DeepSeek and Kimi tracks (90\% equivalence intervals contained in $\pm10$ points). An operation-matched retrieval treatment R runs the same frozen temporal operation but materializes its output in retrieved context rather than callable skill state. Across 18 pre-frozen DeepSeek load-bearing endpoints, T$-$R is exactly 0.0 points with a 90\% endpoint-bootstrap interval of $[-5.6,5.6]$; three held-out Kimi alignment endpoints are also identical under T and R across five fresh repeats. Thus the evidence supports conditional \emph{operation-level} temporal bottlenecks in non-ceiling regimes, but not a material advantage for the callable skill container. The strongest held-out cutoff cell still reaches 100\% under T versus 47.5\% under N, while cross-domain grounding and a fresh Kimi grounding recheck remain null. The resulting claim is narrower and better identified: some recurring temporal operations are binding repairs, and retrieval/context construction can implement the same repair at the tested resolution.
\end{abstract}
